\DSource{0138}{49}
\hypertarget{AB01-PDF0138-P01}{}
\DLine{AB01-PDF0138-body-L001}{}{}{festus est in Luna, ob eius festinum cursum, ita ut fere dimidiae horae sit; nam interdum}
\DLine{AB01-PDF0138-body-L002}{}{}{Lunae motus eo temporis spatio 18 minuta complectitur, cuius duplum est differentia inter}
\DLine{AB01-PDF0138-body-L003}{}{}{maxima et minima nychthemera.}
\hypertarget{AB01-PDF0138-P02}{}
\DLine{AB01-PDF0138-body-L004}{}{}{\Indent Haec differentia e duobus elementis constat, quorum alterum est inaequalitas motus Solis,}
\DLine{AB01-PDF0138-body-L005}{5}{}{sive aequatio, alterum differentia transitus signorum per medium Caelum, quia non omnia illic}
\DLine{AB01-PDF0138-body-L006}{}{}{eadem quantitate ascendunt. Quod maximi ex anomalia motus Solis provenit est fere $3^\circ\qfrac{1}{4}+$}
\DLine{AB01-PDF0138-body-L007}{}{}{$\qfrac{1}{10}$ (1) [$=3^\circ21'$], quod e culminatione signorum $4^\circ\qfrac{1}{4}+\qfrac{1}{5}$ (2) [$=4^\circ27'$] circiter; summa}
\DLine{AB01-PDF0138-body-L008}{}{}{igitur est $7^\circ48'$ (3), quae $\qfrac{1}{2}+\qfrac{1}{50}$ horae aequinoctialis [$=0^h\,31^m\,12^s$] circiter efficiunt. Locus}
\DLine{AB01-PDF0138-body-L009}{}{}{deminutionis fere a $\qfrac{2}{3}$ (4) Aquarii ad initium circiter Scorpionis porrigitur; locus augmenti}
\DLine{AB01-PDF0138-body-L010}{10}{}{ab initio Scorpionis ad $\qfrac{2}{3}$ (5) fere Aquarii. --- In tabulis huius nostri libri nos motus medios}
\DLine{AB01-PDF0138-body-L011}{}{}{descripsimus, ponentes locum Solis, iuxta eius motum medium, in $18^\circ19'$ [Aquarii], at secun-}
\DLine{AB01-PDF0138-body-L012}{}{}{dum veracem motum apparentem in $20^\circ$ Aquarii. Et iuxta hoc nychthemerum reliqua nych-}
\DLine{AB01-PDF0138-body-L013}{}{p. 75.}{themera anni in hoc libro supputa ({\itshape b}).}
\hypertarget{AB01-PDF0138-P03}{}
\DLine{AB01-PDF0138-body-L014}{}{}{\Indent Dicit [auctor]: Si vis dies inaequales convertere in dies medios, ad quos in tabulis sunt}
\DLine{AB01-PDF0138-body-L015}{15}{}{motus medii descripti, sume gradus inter medium locum Solis datum et locum alterum ad}
\DLine{AB01-PDF0138-body-L016}{}{}{quem Sol motu medio pervenerit ({\itshape c}); item, per tempora ascensionum rectarum signorum,}
\DLine{AB01-PDF0138-body-L017}{}{}{gradus inter primum locum verum et locum ad quem Sol motu vero pervenerit. Si numerus}
\DLine{AB01-PDF0138-body-L018}{}{}{horum temporum numerum graduum motus medii superat, quanta pars unius horae aequino-}
\DLine{AB01-PDF0138-body-L019}{}{}{ctialis differentia sit, vide, et eam partem diebus inaequalibus datis adde. Si contra numerus}
\DLine{AB01-PDF0138-body-L020}{20}{}{temporum minor est, eam partem de iis diebus deme. Exibunt dies medii, in quos dies inaequa-}
\DLine{AB01-PDF0138-body-L021}{}{}{les dati convertendi erant, sive a meridie sive a media nocte calculus diei inchoatus est. ---}
\DLine{AB01-PDF0138-body-L022}{}{}{Rem contrariam fac, si dies medios tabularum in dies inaequales convertere vis; idest partem}
\DLine{AB01-PDF0138-body-L023}{}{}{illam diebus mediis adde vel de iis deme, prout numerus temporum minor vel maior fuerit.}
\DLine{AB01-PDF0138-body-L024}{}{}{Invenies dies inaequales quaesitos.}
\hypertarget{AB01-PDF0138-P04}{}
\DLine{AB01-PDF0138-body-L025}{25}{}{\Indent Iuxta hoc principium, in hoc nostro libro constitutum, numerus temporum semper minor}
\DLine{AB01-PDF0138-body-L026}{}{}{erit a loco Solis dato usque ad terminum longi spatii temporis, in quo variatio apogei Solis,}
\DLine{AB01-PDF0138-body-L027}{}{}{quod in ecliptica invenimus, crescet; qua re mutabitur quod ex inaequalitate motus Solis pro-}
\DLine{AB01-PDF0138-body-L028}{}{}{venit. Et cum res ita sint ut diximus, loco medio Lunae in initio calculi 18 (6) minuta addidi-}
\DLine{AB01-PDF0138-body-L029}{}{}{mus; portionem autem variationis nychthemerōn, quae ad singulos gradus signorum pertinet,}
\DLine{AB01-PDF0138-body-L030}{30}{}{in tabulis ascensionum rectarum (7) descripsimus in columna quae ascensiones singulorum gra-}
\DLine{AB01-PDF0138-body-L031}{}{}{duum sequitur. Si igitur quod vero gradui Solis respondet sumpserimus, et partem propor-}
\DLine{AB01-PDF0138-body-L032}{}{}{tionalem unius horae aequinoctialis de diebus inaequalibus detraxerimus, remanebunt dies}
\DLine{AB01-PDF0138-body-L033}{}{}{medii, ad quos in tabulis motus reperiuntur; si vero eam diebus mediis addiderimus, exibunt}
\DLine{AB01-PDF0138-body-L034}{}{}{dies inaequales qui calculo deducuntur ({\itshape d}).}
\par\vspace{6pt}\hrule height.25pt\vspace{5pt}
\noindent\begin{minipage}[t]{.48\textwidth}\vspace{0pt}\fontsize{9}{11.5}\selectfont\setlength{\GutterOffset}{0pt}
\hypertarget{AB01-PDF0138-N01}{}
\DLine{AB01-PDF0138-notes_left-L001}{}{}{\Indent (1) Male Plato $3^\circ\qfrac{1}{5}+\qfrac{1}{10}$, quod $3^\circ18'$ effice-}
\DLine{AB01-PDF0138-notes_left-L002}{}{}{ret. --- In {\itshape Almag.} III, 8 (ed. Halma t. I, 209) est $3^\circ\qfrac{2}{3}$;}
\DLine{AB01-PDF0138-notes_left-L003}{}{}{ibi enim excentricitas orbis solaris maior est quam}
\DLine{AB01-PDF0138-notes_left-L004}{}{}{apud al-Battānī.}
\hypertarget{AB01-PDF0138-N02}{}
\DLine{AB01-PDF0138-notes_left-L005}{}{}{\Indent (2) {\itshape Almag.} $4^\circ\qfrac{2}{3}$; differentia e diversa eclipti-}
\DLine{AB01-PDF0138-notes_left-L006}{}{}{cae obliquitate provenit.}
\end{minipage}
\hfill\begin{minipage}[t]{.48\textwidth}\vspace{0pt}\fontsize{9}{11.5}\selectfont\setlength{\GutterOffset}{0pt}
\hypertarget{AB01-PDF0138-N03}{}
\DLine{AB01-PDF0138-notes_right-L001}{}{}{\Indent (3) {\itshape Almag.} $8^\circ\qfrac{1}{3}$ scilicet $\qfrac{1}{2}+\qfrac{1}{18}$ horae.}
\hypertarget{AB01-PDF0138-N04}{}
\DLine{AB01-PDF0138-notes_right-L002}{}{}{\Indent (4) {\itshape Almag.} $\qfrac{1}{2}$.}
\hypertarget{AB01-PDF0138-N05}{}
\DLine{AB01-PDF0138-notes_right-L003}{}{}{\Indent (5) {\itshape Almag.} $\qfrac{1}{2}$.}
\hypertarget{AB01-PDF0138-N06}{}
\DLine{AB01-PDF0138-notes_right-L004}{}{}{\Indent (6) Plato 58. --- Cur id fecerit explicatur in}
\DLine{AB01-PDF0138-notes_right-L005}{}{}{adnot. {\itshape b} ad hoc caput.}
\hypertarget{AB01-PDF0138-N07}{}
\DLine{AB01-PDF0138-notes_right-L006}{}{}{\Indent (7) Fol. 179,r.--181,v.}
\end{minipage}
\par\vspace{5mm}\hfill{\fontsize{8}{10}\selectfont 7}\hspace{10mm}
